\chapter{Chapter09 — Large Fourier Coefficient}
%%% generated by the blueprint pipeline from FABL.Chapter09.LargeFourierCoefficient
\section{Overview}

Declarations of the Lean module \texttt{FABL.Chapter09.LargeFourierCoefficient}.

\section{Declarations}

\begin{theorem}[Friedgut degree at one eighth]
\lean{FABL.friedgutDegree_one_eighth}
\label{FABL.friedgutDegree_one_eighth}
\leanok
\uses{FABL.BooleanFunction, FABL.BooleanFunction.toReal, FABL.friedgutDegree, FABL.totalInfluence}
\difficulty{3}
For a Boolean function $f$ on the $n$-dimensional sign cube, its Friedgut degree at
parameter $1/8$ equals the least integer greater than or equal to eight times its total
influence: $\lceil 8 I(f)\rceil$, where $I(f)$ denotes its total influence.
\end{theorem}

\begin{theorem}[Card friedgut coordinates one eighth le]
\lean{FABL.card_friedgutCoordinates_one_eighth_le}
\statementsource{Analysis-of-Boolean-Functions-by-Ryan-ODonnell}{
  pdf-ffbe9c7b3306-p269-b004,
  pdf-ffbe9c7b3306-p270-b002
}
\label{FABL.card_friedgutCoordinates_one_eighth_le}
\leanok
\uses{FABL.BooleanFunction, FABL.BooleanFunction.toReal, FABL.card_friedgutCoordinates_le, FABL.friedgutCoordinates, FABL.friedgutDegree_one_eighth, FABL.totalInfluence}
\difficulty{4}
% dataset-pipeline: stub -- prose pending build_informal_statements.py
Statement pending regeneration.
\end{theorem}

\begin{definition}[Low-degree subsets]
\lean{FABL.lowDegreeSubsets}
\statementsource{Analysis-of-Boolean-Functions-by-Ryan-ODonnell}{
  pdf-ffbe9c7b3306-p028-b002,
  pdf-ffbe9c7b3306-p270-b002
}
\label{FABL.lowDegreeSubsets}
\leanok
Given a finite coordinate set $J\subseteq\{1,\ldots,n\}$ and $k\in\mathbb N$, the
low-degree subsets of $J$ are all subsets of $J$ having cardinality at most $k$.
\end{definition}

\begin{theorem}[Mem low degree subsets]
\lean{FABL.mem_lowDegreeSubsets}
\statementsource{Analysis-of-Boolean-Functions-by-Ryan-ODonnell}{pdf-ffbe9c7b3306-p270-b002}
\label{FABL.mem_lowDegreeSubsets}
\leanok
\uses{FABL.lowDegreeSubsets}
\difficulty{1}
% dataset-pipeline: stub -- prose pending build_informal_statements.py
Statement pending regeneration.
\end{theorem}

\begin{theorem}[Counting subsets of bounded size]
\lean{FABL.card_lowDegreeSubsets_eq_sum_choose}
\label{FABL.card_lowDegreeSubsets_eq_sum_choose}
\leanok
\uses{FABL.lowDegreeSubsets}
\difficulty{4}
Let $J$ be a subset of a finite set and let $k$ be a nonnegative integer. The number of
subsets of $J$ having size at most $k$ is $\displaystyle\sum_{j=0}^{k}\binom{|J|}{j}$.
\end{theorem}

\begin{theorem}[Upper bound for low degree subsets]
\lean{FABL.card_lowDegreeSubsets_le}
\label{FABL.card_lowDegreeSubsets_le}
\leanok
\uses{FABL.card_lowDegreeSubsets_eq_sum_choose, FABL.lowDegreeSubsets}
\difficulty{4}
Let $J$ be a subset of a finite set and let $k$ be a nonnegative integer. The number of
subsets of $J$ having size at most $k$ is at most $(k+1)(|J|+1)^k$.
\end{theorem}

\begin{theorem}[Parseval identity for a finite family]
\lean{FABL.sum_fourierWeight_mem_add_outside_eq_one}
\proofsource{Analysis-of-Boolean-Functions-by-Ryan-ODonnell}{
  pdf-ffbe9c7b3306-p025-b006,
  pdf-ffbe9c7b3306-p027-b007
}
\label{FABL.sum_fourierWeight_mem_add_outside_eq_one}
\leanok
\uses{FABL.BooleanFunction, FABL.BooleanFunction.toReal, FABL.fourierCoeff, FABL.fourierWeight, FABL.fourierWeightOutside, FABL.sum_sq_fourierCoeff_eq_one}
\difficulty{4}
Let $n$ be a natural number, let $f:\{-1,1\}^n\to\{-1,1\}$, and let $\mathcal U$ be a
finite family of subsets of the coordinate set $\{1,\ldots,n\}$. For each such subset
$S$, write $\widehat f(S)=\mathbb E_x\bigl[f(x)\prod_{i\in S}x_i\bigr]$, where $x$ is
uniform on the sign cube. Then $\sum_{S\in\mathcal U}\widehat
f(S)^2+\sum_{S\notin\mathcal U}\widehat f(S)^2=1$, with the second sum over all subsets
of the coordinate set.
\end{theorem}

\begin{definition}[Nonempty low-degree coordinate family]
\lean{FABL.corollary9_32Family}
\statementsource{Analysis-of-Boolean-Functions-by-Ryan-ODonnell}{pdf-ffbe9c7b3306-p270-b002}
\label{FABL.corollary9_32Family}
\leanok
\uses{FABL.BooleanFunction, FABL.friedgutCoordinates, FABL.friedgutDegree, FABL.lowDegreeSubsets}
For a Boolean function $f$ on the $n$-dimensional sign cube, take the subsets of size at
most the Friedgut degree of $f$ at parameter $1/8$ from the coordinate set selected for
$f$ at that parameter, and remove the empty subset.
\end{definition}

\begin{theorem}[Membership in the nonempty low-degree family]
\lean{FABL.mem_corollary9_32Family}
\label{FABL.mem_corollary9_32Family}
\leanok
\uses{FABL.BooleanFunction, FABL.corollary9_32Family, FABL.friedgutCoordinates, FABL.friedgutDegree, FABL.mem_lowDegreeSubsets}
\difficulty{3}
Let $f$ be a Boolean function on the $n$-dimensional sign cube, and let $S$ be a subset
of its coordinates. Then $S$ belongs to the family associated with $f$ exactly when $S$
is a nonempty subset of the Friedgut coordinates of $f$ at parameter $1/8$, and its
cardinality is at most the Friedgut degree of $f$ at parameter $1/8$.
\end{theorem}

\begin{theorem}[Fourier mass on a finite family]
\lean{FABL.one_fourth_le_sum_corollary9_32Family}
\proofsource{Analysis-of-Boolean-Functions-by-Ryan-ODonnell}{pdf-ffbe9c7b3306-p270-b002}
\label{FABL.one_fourth_le_sum_corollary9_32Family}
\leanok
\uses{FABL.BooleanFunction, FABL.BooleanFunction.toReal, FABL.corollary9_32Family, FABL.fourierCoeff, FABL.fourierWeightOutside, FABL.friedgutCoordinates, FABL.friedgutDegree, FABL.friedgut_spectrum_concentrated, FABL.lowDegreeSubsets, FABL.mean_eq_fourierCoeff_empty, FABL.sum_fourierWeight_mem_add_outside_eq_one, FABL.totalInfluence, FABL.variance, FABL.variance_eq_four_mul_probabilities, FABL.variance_le_totalInfluence}
\difficulty{5}
Let $n\in\mathbb{N}$ and let $f:\{-1,1\}^n\to\{-1,1\}$. If the variance of $f$ under the
uniform distribution on the sign cube is at least $1/2$, then the sum of the squares of
its Fourier coefficients over the family consisting of the nonempty subsets of the
Friedgut coordinates of $f$ at accuracy $1/8$ having size at most the corresponding
Friedgut degree is at least $1/4$.
\end{theorem}

\begin{theorem}[Large low-degree Fourier coefficient]
\lean{FABL.exists_large_lowDegree_fourierCoeff}
\proofsource{Analysis-of-Boolean-Functions-by-Ryan-ODonnell}{
  pdf-ffbe9c7b3306-p270-b001,
  pdf-ffbe9c7b3306-p270-b002
}
\label{FABL.exists_large_lowDegree_fourierCoeff}
\leanok
\uses{FABL.BooleanFunction, FABL.BooleanFunction.toReal, FABL.card_lowDegreeSubsets_le, FABL.corollary9_32Family, FABL.fourierCoeff, FABL.friedgutCoordinates, FABL.friedgutDegree, FABL.mem_corollary9_32Family, FABL.one_fourth_le_sum_corollary9_32Family, FABL.variance}
\difficulty{5}
Let $n$ be a nonnegative integer and let $f:\{-1,1\}^n\to\{\pm1\}$ have variance at
least $1/2$ when viewed as a real-valued function on the sign cube. Write $d$ for its
Friedgut degree at parameter $1/8$ and $J$ for its Friedgut coordinate set at that
parameter. Then some nonempty set $S\subseteq\{0,\ldots,n-1\}$ satisfies $|S|\le d$ and
$\frac{1}{4(d+1)(|J|+1)^d}\le \widehat f(S)^2$, where $\widehat f(S)$ is the Fourier
coefficient of the real-valued form of $f$ at $S$.
\end{theorem}

\begin{theorem}[Variance and the constant Fourier coefficient]
\lean{FABL.one_half_le_fourierCoeff_empty_sq_of_variance_le}
\proofsource{Analysis-of-Boolean-Functions-by-Ryan-ODonnell}{pdf-ffbe9c7b3306-p270-b003}
\label{FABL.one_half_le_fourierCoeff_empty_sq_of_variance_le}
\leanok
\uses{FABL.BooleanFunction, FABL.BooleanFunction.toReal, FABL.fourierCoeff, FABL.mean_eq_fourierCoeff_empty, FABL.variance, FABL.variance_eq_four_mul_probabilities}
\difficulty{3}
Let $n\in\mathbb{N}$ and let $f:\{-1,1\}^n\to\{-1,1\}$. If the uniform variance of $f$
is at most $1/2$, then the square of its Fourier coefficient at the empty set is at
least $1/2$.
\end{theorem}
